掷一枚不均匀硬币，正面概率 0.3。定义 \(X(\text{正面})=1\)，\(X(\text{反面})=0\)，则 \(P(X=1)=0.3\)，\(P(X=0)=0.7\)。\(X\) 是样本空间 \(\{\text{正面},\text{反面}\}\) 上的函数——随机性来自输入结果，函数规则本身是确定的。

本单元按统一概念到具体表示的顺序展开：

\begin{enumerate}
\tightlist
\item
  随机变量是样本空间上的函数：理解随机性来自输入结果，而变量本身是一个确定函数。
\item
  离散分布与概率质量函数：用 PMF 与 CDF 描述可数取值。
\item
  连续分布、密度与分位数：理解密度不是概率，CDF 才是离散与连续的统一对象。
\item
  随机变量变换与逆变换抽样：当 \(Y=g(X)\) 时，怎样从 \(X\) 的分布得到 \(Y\) 的分布。
\end{enumerate}

随后\hyperref[ch:040]{``期望与矩：怎样用少量数字概括一个分布''}会进一步把整个分布压缩成位置、波动与关联的数值摘要。

\subsection{一个具体算例}

\(X\sim\Bernoulli(0.3)\)：PMF 为 \(p(1)=0.3\), \(p(0)=0.7\)；CDF 为 \(F(x)=0\)（\(x<0\)），\(F(x)=0.7\)（\(0\le x<1\)），\(F(x)=1\)（\(x\ge 1\)）。令 \(Y=2X+1\)，则 \(P(Y=1)=0.7\)，\(P(Y=3)=0.3\)。

连续情形：\(X\sim\Uniform(0,1)\)，密度 \(f(x)=\mathbf{1}_{[0,1]}(x)\)，CDF \(F(x)=x\)（\(x\in[0,1]\)）。令 \(Y=-\ln X\)，则 \(F_Y(y)=P(-\ln X\le y)=P(X\ge e^{-y})=1-e^{-y}\)（\(y>0\)），即 \(Y\sim\Exp(1)\)——这是逆变换抽样的一个实例。

\subsection{怎么读这一章}

四篇可以按所需表示进入：

\begin{itemize}
\tightlist
\item
  若仍把随机变量理解成会随机变化的字母，先读\hyperref[sec:05-Probability/03-Random-Variables/01]{``随机变量是样本空间上的函数''}；随机性来自输入样本点，函数规则本身是确定的。
\item
  若取值有限或可数，读\hyperref[sec:05-Probability/03-Random-Variables/02]{``离散分布与概率质量函数''}；PMF 把概率直接放在点上，CDF 则累计这些质量。
\item
  若通过区间积分计算概率，读\hyperref[sec:05-Probability/03-Random-Variables/03]{``连续分布、密度与分位数''}；密度是单位长度上的局部浓度，不是单点概率，而且并非每个分布都拥有 PDF。
\item
  若定义了 \(Y=g(X)\) 或需要从均匀随机数生成目标分布，读\hyperref[sec:05-Probability/03-Random-Variables/04]{``随机变量变换与逆变换抽样''}；先把 \(Y\) 的事件拉回 \(X\) 空间，再处理尺度伸缩与多个原像。
\end{itemize}

阅读时在能区分四层对象时停下：样本点 \(\omega\)、随机变量函数 \(X\)、实现值 \(x\) 与分布 \(P_X\)。同分布只说明最后一层相同，不说明随机变量在同一实验中相等或具有相同依赖关系。

\subsection{本章篇目}

以下篇目按概念依赖排列；若从具体问题进入，也可以先打开最接近当前困惑的一篇，再沿篇内链接回补。

\begin{itemize}
\tightlist
\item \hyperref[sec:05-Probability/03-Random-Variables/01]{``随机变量是样本空间上的函数''}；
\item \hyperref[sec:05-Probability/03-Random-Variables/02]{``离散分布与概率质量函数''}；
\item \hyperref[sec:05-Probability/03-Random-Variables/03]{``连续分布、密度与分位数''}；
\item \hyperref[sec:05-Probability/03-Random-Variables/04]{``随机变量变换与逆变换抽样''}。
\end{itemize}
